\section{The instrument: the r2m3 plant, its V\&V record and its numbered limits}\label{sec:the_instrument}

This section describes the plant as built and as read back from the simulator, and states, requirement by requirement, what supports the comparison in Section~\ref{sec:results} and which limits every result carries.
The plant is called r2m3 (two on-ramps, three main lanes) and its frozen configuration is \texttt{plant\_c3}.
The standard we apply is ``valid for its intended use''; we do not call the plant proven.
Grey bracketed codes such as \ledger{V367} point to entries of the project's decision ledger (Appendix~\ref{sec:supplement}).
\todo{release form of the ledger and code is an open decision (plan.md, open decision on release)}

\subsection{Intended use}\label{sec:intended_use}

The V\&V record fixes the use before it lists any requirement:

\begin{quote}
\emph{Use:} compare controllers that set the time headway of automated cars (ALKS; 1.6, 1.2, 1.0 or 0.8\,s) against rule-based controllers and no control, on paired differences in door-to-door delay (with insertion queues), served flow, close-following (TTC below 3\,s), braking at step resolution and severe events, on a synthetic three-lane freeway with two on-ramps, 50\,\% ALKS, demand 4,090--6,000\,veh/h, steady or with tapered slowdowns (X1, X2).

\emph{Not for:} absolute travel times or capacities of a real road; effects on a calibrated or field site; lateral or merge-conflict safety; V2V with latency or loss; human adaptation to automated cars.

\emph{Acceptance:} the plant is valid for this use when every requirement below is met, or its gap is stated as a limit that every claim sentence carries.
\end{quote}

ALKS denotes the Automated Lane Keeping System of UN Regulation No.~157~\autocite{UNR157}\verify; HDV denotes a human-driven vehicle.
Some reported rows fall outside this use, and we label them ``outside the plant's stated intended use'' instead of widening it \ledger{V462}.
They are the demand band and the training condition T of one opponent, both under a random training slowdown (one slowdown on edge 3, 4 or 5 at a limit drawn from 25--60\,km/h), and the two history cells H1 (a slow truck) and H2 (ramp bursts).
The band's demand, 4,500--6,000\,veh/h, lies inside the stated range; only its slowdown does not.
History rows from the earlier plant versions P0 and P0d are labelled ``not the V\&V target''.
The intended use was written on 2026-11-07, after confirmatory~3 and round~1 had already used the plant \ledger{V459}; Section~\ref{sec:verification_changed} returns to this.

\subsection{Network}\label{sec:network}

The network (Fig.~\ref{fig:network}, Table~\ref{tab:plant}) is a 6.8-km freeway with three main lanes and two single-lane on-ramps, flat, with no off-ramp and no lane drop.
Every lane, ramps included, has a 100-km/h limit and is 3.2\,m wide.
Each ramp joins through an acceleration lane (3a and 4a), 294.72\,m long against a nominal 300\,m (L1).
The merges are ``1.2\,km apart'' measured from the end of 3a to the start of 4a (edge~3, 1,194.3\,m); junction to junction they are 1,500\,m apart, and we state both.
Ramp~1 carries a signal that is held green under every controller in this paper; ramp~2 has neither signal nor detectors.
The three route lengths, about 6,778, 3,700 and 2,211\,m without junction-internal lanes, are derived from the network file, not read from a specification.

\begin{figure}[tbp]
\centering
\todo{F1: r2m3 to scale, edges, lanes, acceleration lanes, ramps, X1 and X2 edges, detectors, both reference distances}
\caption{The r2m3 network to scale. X2 acts on edge~3 between the merges and X1 on edge~4 downstream of 4a.}\label{fig:network}
\end{figure}

\begin{table}[tbp]
\centering
\todo{T2: geometry, demand, fleet (configured and realised), actuator, slowdowns, simulation settings}
\caption{The plant at a glance. Route lengths are derived. The SUMO arguments are those stored with the r1m3 baseline.}\label{tab:plant}
\end{table}

\subsection{Demand, run window and flow}\label{sec:demand}

Vehicles enter as one Poisson stream at $q$\,veh/h, constant within a run, and split over the three routes in shares 70/15/15 (mainline, ramp~1, ramp~2), departing at the highest speed SUMO allows on insertion (\texttt{departSpeed} max).
A run lasts 1,800\,s: a 300-s warm-up and a 1,500-s measurement window.
Served flow is the number of vehicles that reach edge~4 inside the window, written as vehicles per 1,500\,s; for scale, 2,132 vehicles correspond to about 5,120\,veh/h.
No test yet shows that the inserted demand matches the configured level and ramp shares within sampling error (L2).

\subsection{Fleet and drivers}\label{sec:fleet}

Each vehicle's class is drawn on insertion: 7.5\,\% trucks, 50\,\% ALKS cars and 42.5\,\% HDV cars as configured.
A read-back over 106,128 draws gave 42.76\,\% HDV, 7.68\,\% trucks and 49.56\,\% ALKS; these draws used the routes of the single-ramp plant r1m3, not r2m3 (L4).
The full parameter table, with a source and status per parameter, is Table~T11 in the supplement.

HDVs and trucks follow SUMO's Krauss model~\autocite{Krauss1998}\verify.
The HDV time-gap scale (0.8) and action-step scale (0.5) were calibrated earlier on a human-only lane-drop corridor \ledger{V187, V188}, not on this plant.
As a result the realised mean HDV gap parameter is 0.934\,s, against a mean of 1.15\,s for close following measured on a UK motorway~\autocite{AlObaedi2019}\verify\ (L5).
HDV acceleration 2.6, deceleration 4.5 and emergency deceleration 9.0\,m/s$^2$ are SUMO defaults and lie above the ranges in the sources we collected (for acceleration, 1.8--2.2\,m/s$^2$~\autocite{Long2000}\verify); the lane-change assertiveness of 1.4 has no source (L5, L8).

ALKS cars follow the Intelligent Driver Model~\autocite{Treiber2000}\verify\ with a default time headway of 1.6\,s.
That value is the R157 minimum gap at 60\,km/h~\autocite{UNR157}\verify; the regulation's table stops at 60\,km/h, while the plant runs ALKS at up to 100\,km/h and commands gaps down to 0.8\,s.
Whether these gaps and deceleration limits are consistent with R157 has not been analysed (L7).
Emission classes are read back (cars Euro-4 petrol, trucks Euro-VI 34--40\,t, HBEFA4) and fuel is reported only with that disclosure.
Every parameter source carries [VERIFY] until the author has read it in full.

\subsection{Actuator and sensing}\label{sec:actuator}

Fig.~\ref{fig:boundary} draws the boundary between SUMO, our environment layer and the controllers.
A controller chooses one of four headway levels (1.6, 1.2, 1.0, 0.8\,s) for each ALKS every 10\,s and holds it for 10\,s.
The command takes effect only when the ALKS follows another ALKS within 150\,m and has had the same leader on two consecutive updates; otherwise the ALKS keeps 1.6\,s.
A longer commanded gap is approached at no more than 0.2\,s of headway per second (\texttt{tau\_rise}), so that a command cannot force hard braking; the ALKS action step is 0.5\,s.
V2V messages and roadside speeds are ideal, with no latency or loss (L15).

The environment layer also applies a smooth no-undertaking rule to ALKS, which is our own model and replaces SUMO's abrupt rule for them (L9).
Above 60\,km/h and outside the acceleration lanes, an ALKS's speed is capped at that of a slower vehicle ahead in the lane to its left within 8\,s of travel; the ALKS eases down to the cap at the larger of a speed ratio of 1.1 per second and the deceleration needed to match that vehicle over the gap, the latter limited to 4.4\,m/s$^2$.
No inspection yet shows that the policy inputs carry no information a real vehicle could not have, such as the seed, the event timing or a future state (L14).

\begin{figure}[tbp]
\centering
\todo{F2: SUMO, the environment layer (V2V gate, tau\_rise, no-undertaking rule) and the controllers; observation, actuation and ground-truth measurement paths}
\caption{Plant, environment code and controllers as built. Measurement counters read ground truth; the policy inputs are not yet inspected for information a vehicle could not have (L14).}\label{fig:boundary}
\end{figure}

\subsection{Slowdowns}\label{sec:slowdowns}

Two slowdowns define the decision cells (Fig.~\ref{fig:slowdowns}).
In X1, edge~4's limit drops to 25\,km/h at 600\,s and again at 1,200\,s; each event lasts 120\,s, including a 20-s linear taper, and the limit is then released at once.
In X2, edge~3's limit drops to 30\,km/h from 700\,s for 180\,s, including the same taper, and is then released at once.
The taper is updated on the 0.1-s simulation clock.
It exists because the untapered events produced artificial hard braking: on one seed at 4,600\,veh/h, 57 of 59 vehicle-seconds of braking at 7\,m/s$^2$ or harder fell in the first 10\,s of the events \ledger{V345}.
The X1 and X2 values have no field source, and the record that checks them event by event is not written (L12).

\begin{figure}[tbp]
\centering
\todo{F3: speed limit against time for X1 and X2 with warm-up and window shaded; inset, the training draw (outside the plant's stated intended use)}
\caption{Speed-limit profiles of the slowdowns X1 and X2. The inset shows the random training slowdown, which lies outside the plant's stated intended use.}\label{fig:slowdowns}
\end{figure}

\subsection{Simulation and reproducibility}\label{sec:reproducibility}

The plant runs in SUMO 1.27.1~\autocite{Lopez2018} through libsumo, one simulation per process, with a 0.1-s step, teleporting off, collisions recorded as warnings, and junction collision checks on.
SUMO uses its ballistic position update because the vehicle types set an action step \ledger{V367}; no \texttt{plant\_c3} log line confirms this yet, and the stored SUMO arguments are those of the r1m3 baseline.
Replays reproduce stored rows: 200 base replays at a frozen commit matched 28,555 stored keys \ledger{V398}.
Rows also match across machines: a rented server reproduced the laptop's check rows and training fingerprint \ledger{V437}, and the laptop reproduced the server's tuning rows on three seeds per cell \ledger{V461}.
Stored result files are hashed and the toolchain is pinned.
Two deviations are disclosed (L17): the metadata of the CMA-ES tuning run records the hash of its script but no commit \ledger{V461}, and round~2 resumes interrupted work by matching rows on name, demand level and seed, without a configuration hash \ledger{V455}.

\subsection{The V\&V record}\label{sec:vv_record}

The record follows the usual systems-engineering pattern: intended use, requirements, a verification method per requirement (inspection, analysis, demonstration or test), evidence, status, traceability to claim elements, and stated limits.
Background sources for this structure are IEEE~1012, Sargent's validation framework and the FHWA Traffic Analysis Toolbox volume~III\verify.
\todo{add the three to references.bib once read}
Its 28 requirements fall in six groups: road, demand and insertion; drivers and fleet; actuator and sensing; numerics and reproducibility; measurement; and behaviour that makes control meaningful.
Table~\ref{tab:vv} summarises them and Table~T10 in the supplement lists them in full.
At the time of writing, 6 are met, 13 partial, 7 open gaps and 2 stated limits.
\todo{regenerate the tally from the reviewed record}

Two pieces of evidence that do not depend on any learned policy support the requirements that the command binds and that control matters (PR-10, PR-25).
In a pre-run on 10 seeds at 5,300\,veh/h, the fixed 0.8-s rule had less finishers' delay than no control by 29.6\,\% [24.8, 48.7] in steady traffic, 19.9\,\% [15.8, 22.1] in X1 and 12.8\,\% [9.0, 24.8] in X2, and more flow by 3.5\,\% [1.6, 4.6], 4.4\,\% [3.0, 5.8] and 3.0\,\% [2.3, 4.8] (95\,\% CIs) \ledger{V382}.
No arm collided, and braking counted at 1\,Hz was 0.965--0.991 of braking counted every 0.1\,s.
This is 10 seeds at one level, so PR-25 stays partial; the record's earlier citation of policy results for PR-25 was circular and is not used (L21).
A steady-traffic check on the earlier plant P0d (r1m3, 20 seeds, four levels; not the V\&V target) passed six of seven criteria, with the seventh counted as met with disclosure, and had no collision or teleport in 80 runs \ledger{V376}.
On \texttt{plant\_c3} itself, no control collided once in 40 X1 runs of round~1 and never in X2, so the collision requirement is partial: collisions are counted and reported for every arm, but are too rare to compare (L16).

\begin{table}[tbp]
\centering
\todo{T3: requirements by group and status, traceability to claim elements, pointer to L1--L27, the three gaps that block ``verified and validated''}
\caption{V\&V summary for r2m3 \texttt{plant\_c3}.}\label{tab:vv}
\end{table}

\subsection{What verification changed, and when}\label{sec:verification_changed}

The plant changed twice, each time after a check found a measured defect (Table~\ref{tab:lineage}).
From P0 to P0d \ledger{V367}: on two seeds at 5,300\,veh/h, ALKS braked at 7\,m/s$^2$ or harder for 1,749 and 972 vehicle-seconds, which the stored 1-Hz metric reported as 0 and 0; the smooth rule's look-ahead and braking limit were then set.
From P0d to \texttt{plant\_c3}: on one seed at 5,300\,veh/h, an instant rise of the commanded gap produced 71.3 and 75.3 vehicle-seconds of braking at 7\,m/s$^2$ in X1 and X2, against 11 and 5 seen at 1\,Hz \ledger{V377}, which led to \texttt{tau\_rise}; the untapered slowdowns described above led to the taper \ledger{V345}.
On two development seeds, and for the fixed 0.8-s rule only, the fixes removed moving full-emergency events (72, 63 and 78 to 0, 0 and 0 in steady traffic, X1 and X2) at about 10--13\,\% more delay under the slowdowns \ledger{V380}; this is a description, not a check.

The order matters for how far the comparison can be trusted.
P0d came after phase~1 and confirmatory~2, and its look-ahead and braking limit were chosen after a check failed on undertaking.
\texttt{plant\_c3} came after confirmatory~2 and a later single test; one rule family was derived from P0 policies, and the reward was shaped by runs that knew the fixed rule.
The V\&V target and the intended use were fixed after confirmatory~3 and round~1.
Each change corrected a defect in what the plant did, not a ranking of controllers, and \texttt{plant\_c3} stays frozen for all later work.
An earlier entry of the ledger asked that controllability properties be fixed before control results \ledger{V245}; for this plant that was not done in that order, and we report it rather than repair it.

\begin{table}[tbp]
\centering
\todo{T4: P0, P0d, plant\_c3; trigger, change, measured effect, timeline against the control results}
\caption{Plant lineage: what verification found, what it changed, and when relative to the control results.}\label{tab:lineage}
\end{table}

\subsection{Limits carried by every claim}\label{sec:limits}

Every requirement that is partial, an open gap or a stated limit becomes a numbered limit, as do two met requirements whose status carries a condition (the no-undertaking rule and the input normalisation) and three record items with no requirement (an uncalibrated model, a synthetic geometry, and no measure of lateral or merge conflicts).
This gives L1--L27 (Table~T19 in the supplement), and every claim sentence in this paper reads ``within limits L1--L27''.
Three of them must close before the words ``verified and validated'' appear anywhere: whether the inputs carry information a vehicle could not have (L14), a traced sample for every reported counter (L20), and the location of gains at merge~2, where 608.5 lane-end stops per run occur against 7 at ramp~1 under no control on P0d, and where 46--61\,\% of an earlier X2 gain sat \ledger{V396} (L24).
When a limit closes, the numbering is regenerated from the reviewed record.
